\section{\texorpdfstring{هفته‌ی 6: ارزیابی مدل‌ها 1: نیکویی برازش،
\lr{overfitting} و
\lr{cross-validation}}{هفته‌ی 6: ارزیابی مدل‌ها 1: نیکویی برازش، overfitting و cross-validation}}

\begin{quote}
\textbf{فصل 3: ارزیابی مدل‌های شناختی} · درس 1 ساعت · \lr{lab} 2 ساعت
\textbf{داده‌ی اصلی:} داده‌های فراموشی هفته‌های 2\lr{--}5، این بار با چهار مدل
رقیب: نمایی (\(k=2\))، توانی (\(k=2\))، نمایی با مجانب (\(k=3\)) و دوردی
(\(k=4\)).
\end{quote}

\begin{center}\rule{0.5\linewidth}{0.5pt}\end{center}

\subsection{واژه‌های کلیدی}

واژه‌های مهم این درس با معادل فارسی‌شان (حتی اگر در درس‌های قبل آمده
باشند).

\begin{longtable}[]{@{}
  >{\raggedright\arraybackslash}p{(\linewidth - 4\tabcolsep) * \real{0.3333}}
  >{\raggedright\arraybackslash}p{(\linewidth - 4\tabcolsep) * \real{0.3333}}
  >{\raggedright\arraybackslash}p{(\linewidth - 4\tabcolsep) * \real{0.3333}}@{}}
\caption{\centering واژه‌های کلیدی}\tabularnewline
\toprule\noalign{}
\begin{minipage}[b]{\linewidth}\raggedright
\lr{English}
\end{minipage} & \begin{minipage}[b]{\linewidth}\raggedright
فارسی
\end{minipage} & \begin{minipage}[b]{\linewidth}\raggedright
توضیح کوتاه
\end{minipage} \\
\midrule\noalign{}
\endfirsthead
\toprule\noalign{}
\begin{minipage}[b]{\linewidth}\raggedright
\lr{English}
\end{minipage} & \begin{minipage}[b]{\linewidth}\raggedright
فارسی
\end{minipage} & \begin{minipage}[b]{\linewidth}\raggedright
توضیح کوتاه
\end{minipage} \\
\midrule\noalign{}
\endhead
\bottomrule\noalign{}
\endlastfoot
\lr{goodness of fit} & نیکویی برازش & اینکه مدل داده‌ی موجود را چقدر خوب
توصیف می‌کند \\
\lr{fit} / \lr{fitting} & برازش & تنظیم پارامترها تا مدل با داده جور
شود \\
\lr{overfitting} & بیش‌برازش \lr{(overfitting)} & برازش نویز نمونه به‌جای
الگوی واقعی \\
\lr{forgetting curve} & منحنی فراموشی & کاهش یادآوری با گذشت زمان \\
\lr{exponential} / \lr{power function} & تابع نمایی / توانی & دو شکل
رایج برای منحنی فراموشی و یادگیری \\
\lr{trace plot} & نمودار \lr{trace} & مقدار نمونه‌ها در برابر شماره‌ی
تکرار \\
\lr{model flexibility} / \lr{complexity} & انعطاف‌پذیری / پیچیدگی مدل &
دامنه‌ی الگوهایی که مدل می‌تواند تولید کند \\
\lr{training} / \lr{test data} & داده‌ی آموزش / آزمون & داده‌ی برازش /
داده‌ی کنارگذاشته برای سنجش \\
\lr{cross-validation (CV)} & اعتبارسنجی متقابل \lr{(cross-validation)} &
برازش روی بخشی از داده و آزمون روی بخش کنارگذاشته \\
\lr{log-likelihood} / \lr{NLL} & لگاریتم درست‌نمایی / منفی آن & مقداری که
بهینه‌سازها کمینه یا بیشینه می‌کنند \\
\lr{likelihood} & درست‌نمایی & احتمال داده به‌عنوان تابعی از پارامترها \\
\lr{generalization} & تعمیم & پیش‌بینی خوب داده‌ی جدید \\
\lr{parameter} & پارامتر & کمیت قابل تنظیم مدل که بین افراد یا شرایط فرق
می‌کند \\
\lr{SSE} / \lr{RMSD} & مجموع مربعات خطا / ریشه‌ی میانگین مربع انحراف &
معیارهای فاصله‌ی داده و پیش‌بینی \\
\lr{noise} & نویز & تغییرات تصادفی در شواهد یا در داده \\
\lr{response bias} & سوگیری پاسخ & تمایل به یک پاسخ، مستقل از محرک \\
\lr{criterion (c)} & ملاک & آستانه‌ی تصمیم؛ نشان‌دهنده‌ی سوگیری پاسخ \\
\lr{bias\lr{--}variance trade-off} & بده‌بستان سوگیری و واریانس & مدل ساده
سوگیر، مدل منعطف پرواریانس \\
\lr{estimate} / \lr{estimator} & برآورد / برآوردگر & مقدار به‌دست‌آمده
برای پارامتر / روش به دست آوردن آن \\
\lr{data leakage} & نشت داده & رسیدن اطلاعات داده‌ی آزمون به برازش \\
\lr{AIC} / \lr{BIC} & معیار اطلاعاتی آکائیکه / بیزی & برازش به‌علاوه‌ی
جریمه برای تعداد پارامتر \\
\lr{benchmark} & \lr{benchmark} (مجموعه‌آزمون استاندارد) & مجموعه سؤال
برای مقایسه‌ی مدل‌های هوش مصنوعی \\
\lr{benchmark contamination} & آلودگی \lr{benchmark} & بودن سؤال‌های
آزمون در داده‌ی آموزش مدل \\
\lr{parameter recovery} & بازیابی پارامتر & آیا برازش، پارامترهای معلوم
شبیه‌سازی را برمی‌گرداند \\
\end{longtable}

\begin{center}\rule{0.5\linewidth}{0.5pt}\end{center}

\subsection{اهداف یادگیری}

در پایان این هفته، دانشجویان می‌توانند:

\begin{enumerate}
\def\labelenumi{\arabic{enumi}.}
\tightlist
\item
  با استفاده از مفاهیم انعطاف‌پذیری و \lr{overfitting} \textbf{توضیح
  دهند} که چرا برازش خوب شاهد لازم ولی ناکافی برای یک مدل است.
\item
  با شبیه‌سازی \lr{overfitting} را \textbf{نشان دهند}: خطای آموزش با
  پیچیدگی مدل یکنواخت کم می‌شود ولی خطای پیش‌بینی نه.
\item
  اعتبارسنجی \lr{hold-out}، \(K\)\lr{-fold} و \lr{leave-one-out} را با
  امتیاز لگاریتم درست‌نمایی خارج از نمونه \textbf{صورت‌بندی کنند} و
  داده‌های آموزش و آزمون را درست جدا نگه دارند.
\item
  بر اساس ادعای تعمیمی که مطرح است، \textbf{انتخاب کنند} چه چیزی را کنار
  بگذارند \lr{(trial}ها، شرایط، شرکت‌کننده‌ها).
\item
  نتایج \lr{cross-validation}، از جمله عدم قطعیتشان، را برای مقایسه‌ی
  مدل‌های شناختی \textbf{تفسیر کنند}.
\end{enumerate}

\begin{center}\rule{0.5\linewidth}{0.5pt}\end{center}

\subsection{مقدمه}

در تمرین 4 هفته‌ی 3، مدل سه‌پارامتری نمایی با مجانب داده‌های هفته‌ی 2 را
\emph{بهتر} از مدل توانی برازش کرد (\(\mathrm{SSE} \approx 0.020\) در
برابر \(0.026\))، با اینکه داده‌ها از مدل توانی تولید شده بودند. در بخش
\lr{D3} هفته‌ی 3، یک مدل دوردی چهارپارامتری هم همین‌قدر خوب بود. اگر مدل‌ها
را بر اساس برازش انتخاب کنیم، مدل غلط را انتخاب می‌کنیم و تا ابد پارامتر
اضافه می‌کنیم.

این مشکل مخصوص منحنی‌های فراموشی نیست. یک مدل به اندازه‌ی کافی منعطف
می‌تواند هر داده‌ای، از جمله نویز، را برازش کند. \lr{Roberts} و
\lr{Pashler} به‌طور معروفی پرسیدند که یک برازش خوب واقعاً چقدر قانع‌کننده
است (مطالعه‌ی بیشتر را ببینید). جوابشان: نه خیلی، مگر اینکه بدانیم مدل چه
چیزی را \emph{نمی‌توانست} برازش کند.

راه خروج این است که یادمان بیاید \textbf{چرا} مدل برازش می‌کنیم. می‌خواهیم
\emph{فرایندی} را که داده را تولید کرده بگیریم، و مدلی که فرایند را
گرفته باشد باید \textbf{داده‌ی جدید} از همان فرایند را \textbf{پیش‌بینی
کند}. نویز طبق تعریف تکرار نمی‌شود. پس سؤال از «مدل این داده‌ها را چقدر
خوب برازش می‌کند؟» به «مدل داده‌هایی را که ندیده چقدر خوب پیش‌بینی می‌کند؟»
تغییر می‌کند. \lr{cross-validation} مستقیماً با همین داده‌هایی که داریم به
این سؤال جواب می‌دهد.

\paragraph{این موضوع کجا به کار می‌رود؟}

\begin{itemize}
\tightlist
\item
  \textbf{انتخاب بین نظریه‌ها.} هر وقت دو مدل حافظه، تصمیم یا یادگیری
  مقایسه می‌شوند، برازش خوب به‌تنهایی قانع‌کننده نیست؛ داوران مقاله‌ها
  \lr{cross-validation} یا معیاری می‌خواهند که انعطاف را جریمه کند.
\item
  \textbf{پیش‌بینی در روان‌شناسی بالینی و روان‌پزشکی:} مدل‌هایی که عود
  بیماری، پاسخ به درمان یا تشخیص را پیش‌بینی می‌کنند باید روی
  \emph{بیماران جدید} اعتبارسنجی شوند، نه روی بیمارانی که با آن‌ها برازش
  شده‌اند. بسیاری از پیش‌بینی‌کننده‌های منتشرشده دقیقاً در همین آزمون شکست
  می‌خورند.
\item
  \textbf{یادگیری ماشین:} تقسیم \lr{train/validation/test} و
  \lr{cross-validation} با \lr{k} بخش، روش استاندارد انتخاب مدل و
  \lr{hyperparameter}هاست؛ \lr{overfitting} خطر اصلی است.
\item
  \textbf{ارزیابی سیستم‌های هوش مصنوعی:} نمره‌ی \lr{benchmark} فقط وقتی
  معنا دارد که آیتم‌های آزمون در داده‌ی آموزش نبوده باشند (بخش هوش مصنوعی
  و \lr{LLM}ها را ببینید).
\end{itemize}

\begin{center}\rule{0.5\linewidth}{0.5pt}\end{center}

\Needspace{11\baselineskip}

\subsection{جزوه‌ی درس}

\paragraph{به زبان ساده}

مدلی که پیچ‌های کافی داشته باشد می‌تواند از همه‌ی نقطه‌های داده بگذرد، حتی
از نویز تصادفی. این آن را مدل خوبی نمی‌کند: نویز دفعه‌ی بعد فرق می‌کند، پس
مدلی که نویز را کپی کرده داده‌ی \emph{جدید} را بد پیش‌بینی می‌کند. این
\textbf{\lr{overfitting} (بیش‌برازش)} است. راه‌حل این است که مدل‌ها را با
این بسنجیم که داده‌ای را که \textbf{ندیده‌اند} چقدر خوب پیش‌بینی می‌کنند.
\textbf{\lr{cross-validation}} این کار را با همان داده‌ای که داریم انجام
می‌دهد: بخشی از داده را پنهان کن، مدل را روی بقیه برازش کن، بخش پنهان را
پیش‌بینی کن، و تکرار کن تا هر بخش یک بار پنهان شده باشد. مدلی که بهترین
پیش‌بینی را روی داده‌ی پنهان دارد برنده است. پیام اصلی هفته: برازش بهتر به
\emph{همان} داده دلیلی بر مدل بهتر نیست؛ \emph{پیش‌بینی} بهتر داده‌ی جدید
دلیل است.

\textbf{زمان‌بندی پیشنهادی (60 دقیقه):} معیارهای برازش و محدودیت‌هایشان 10
· \lr{overfitting} و انعطاف‌پذیری 15 · منطق تعمیم 5 ·
\lr{cross-validation} 20 · چه چیزی را کنار بگذاریم 10.

\subsubsection{1. معیارهای نیکویی برازش}

\begin{longtable}[]{@{}
  >{\raggedright\arraybackslash}p{(\linewidth - 6\tabcolsep) * \real{0.2500}}
  >{\raggedright\arraybackslash}p{(\linewidth - 6\tabcolsep) * \real{0.2500}}
  >{\raggedright\arraybackslash}p{(\linewidth - 6\tabcolsep) * \real{0.2500}}
  >{\raggedright\arraybackslash}p{(\linewidth - 6\tabcolsep) * \real{0.2500}}@{}}
\caption{\centering 1. معیارهای نیکویی برازش}\tabularnewline
\toprule\noalign{}
\begin{minipage}[b]{\linewidth}\raggedright
معیار
\end{minipage} & \begin{minipage}[b]{\linewidth}\raggedright
تعریف
\end{minipage} & \begin{minipage}[b]{\linewidth}\raggedright
مقیاس
\end{minipage} & \begin{minipage}[b]{\linewidth}\raggedright
توضیح
\end{minipage} \\
\midrule\noalign{}
\endfirsthead
\toprule\noalign{}
\begin{minipage}[b]{\linewidth}\raggedright
معیار
\end{minipage} & \begin{minipage}[b]{\linewidth}\raggedright
تعریف
\end{minipage} & \begin{minipage}[b]{\linewidth}\raggedright
مقیاس
\end{minipage} & \begin{minipage}[b]{\linewidth}\raggedright
توضیح
\end{minipage} \\
\midrule\noalign{}
\endhead
\bottomrule\noalign{}
\endlastfoot
\lr{RMSD} & \(\sqrt{\mathrm{SSE}/N}\) & واحد داده & راحت خوانده می‌شود؛
وزن برابر فرض می‌کند \\
\(R^2\) & \(1 - \mathrm{SSE}/\sum_i (y_i - \bar y)^2\) & 0 تا 1 (برای
مدل‌های خطی) & محبوب ولی برای مدل‌های غیرخطی گمراه‌کننده: می‌تواند منفی شود،
انعطاف را پاداش می‌دهد، و تحت سلطه‌ی واریانس بزرگ بین شرایط است (برازش یک
افت یکنوا تقریباً خودکار \(R^2\) بالا می‌دهد) \\
لگاریتم درست‌نمایی \(\ell(\hat{\boldsymbol\theta})\)، \lr{NLL} & هفته‌ی 4
& \lr{nat} & معیار طبیعی برای مدل‌های مبتنی بر درست‌نمایی \\
انحراف \(G^2\) & \(2(\ell_{\text{\lr{sat}}} - \ell)\) & \ensuremath{\approx} \(\chi^2\) & آزمون
مطلق عدم برازش (هفته‌ی 4) \\
\end{longtable}

همه‌ی این‌ها \textbf{کفایت توصیفی} را می‌سنجند: اینکه مدل داده‌هایی را که
رویشان برازش شده چقدر خوب بازتولید می‌کند.

\subsubsection{\texorpdfstring{2. \lr{overfitting} و
انعطاف‌پذیری}{2. overfitting و انعطاف‌پذیری}}

هر مشاهده را به شکل \textbf{علامت + نویز} بنویسید:
\(y_i = g(x_i) + \varepsilon_i\) که در آن \(g\) فرایند واقعی است. مدل
برازش‌شده سعی می‌کند \(y_i\) را بازتولید کند، ولی فقط بخش \(g\) است که در
داده‌ی جدید دوباره ظاهر می‌شود.

\begin{itemize}
\tightlist
\item
  مدلی که \textbf{خیلی ساده} باشد بخشی از \(g\) را از دست می‌دهد: عدم
  برازش نظام‌مند، یا \textbf{کم‌برازش} \lr{(underfitting}، سوگیری بالا).
\item
  مدلی که \textbf{خیلی منعطف} باشد خم می‌شود تا \(\varepsilon_i\) خاص این
  نمونه را دنبال کند: \textbf{\lr{overfitting}} (واریانس بالا). برازشش
  روی نمونه بهتر و پیش‌بینی‌اش برای داده‌ی جدید بدتر می‌شود.
\end{itemize}

با افزایش پیچیدگی، خطای آموزش همیشه کم می‌شود، در حالی که خطای مورد
انتظار پیش‌بینی اول کم و بعد زیاد می‌شود: یک شکل \lr{U} (بخش \lr{A lab).}

\textbf{انعطاف‌پذیری چیزی بیش از تعداد پارامترهاست.} دو مدل با \(k\)
یکسان می‌توانند از نظر انعطاف، یعنی دامنه‌ی الگوهای داده‌ای که می‌توانند
تولید کنند، خیلی فرق داشته باشند. پارامتری که به‌صورت نما یا جابه‌جایی
وارد می‌شود معمولاً انعطاف بیشتری می‌خرد تا پارامتری که فقط منحنی را مقیاس
می‌دهد. مدل دوردی می‌تواند منحنی‌های نمایی، شبه‌توانی و مجانبی را تقلید کند،
و دقیقاً به همین دلیل برازش خوبش چشمگیر نیست. معیارهای اطلاعاتی (هفته‌ی 7)
پارامترها را می‌شمارند؛ \lr{cross-validation} و \lr{Bayes factor} (هفته‌ی
10) به انعطاف واقعی پاسخ می‌دهند.

\begin{quote}
\textbf{پیشرفته: تجزیه‌ی سوگیری\lr{--}واریانس.} برای مربع خطا، خطای مورد
انتظار پیش‌بینی در نقطه‌ی \(x\) (روی مجموعه‌داده‌های جدید و نویز جدید)
این‌طور تجزیه می‌شود:
\[\mathbb{E}\big[(y_{\text{\lr{new}}} - \hat f(x))^2\big] = \underbrace{\sigma^2}_{\text{\lr{irreducible noise}}} + \underbrace{\big(g(x) - \mathbb{E}[\hat f(x)]\big)^2}_{\text{\lr{bias}}^2} + \underbrace{\mathrm{Var}[\hat f(x)]}_{\text{\lr{variance}}}.\]
افزودن پارامتر سوگیری را کم و واریانس را زیاد می‌کند. بهترین مدل برای
پیش‌بینی این دو را متعادل می‌کند، و وقتی داده کم است لزوماً مدل واقعی نیست.
این تجزیه در کتاب عظیمی‌فر و احمدی به‌تفصیل آمده است (مطالعه‌ی بیشتر را
ببینید).
\end{quote}

\paragraph{خواندن ریاضی}

تجزیه‌ی \lr{bias\lr{--}variance} (کادر پیشرفته) میانگین خطای مربعی پیش‌بینی را
به سه بخش تقسیم می‌کند:

\begin{itemize}
\tightlist
\item
  \(\sigma^2\)، \textbf{نویز کاهش‌ناپذیر}: حتی فرایند واقعی هم نمی‌تواند
  بخش تصادفی داده‌ی جدید را پیش‌بینی کند. هیچ مدلی از این پایین‌تر نمی‌رود.
\item
  \textbf{\lr{bias²}} \(= (g(x) - \mathbb{E}[\hat f(x)])^2\): پیش‌بینی
  \emph{میانگین} مدل (میانگین روی مجموعه داده‌های ممکن) چقدر از حقیقت دور
  است. مدل‌های خیلی ساده \lr{bias} بزرگی دارند.
\item
  \textbf{\lr{variance}} \(= \mathrm{Var}[\hat f(x)]\): منحنی برازش‌شده
  از یک مجموعه داده به مجموعه داده‌ی دیگر چقدر بالا و پایین می‌پرد. مدل‌های
  خیلی منعطف \lr{variance} بزرگی دارند، چون دنبال نویز هر مجموعه داده
  می‌دوند.
\end{itemize}

اضافه کردن پارامتر، خطا را از جمله‌ی \lr{bias} به جمله‌ی \lr{variance}
منتقل می‌کند. بهترین پیش‌بینی‌کننده جایی است که جمع این دو کمترین باشد؛ به
همین دلیل وقتی داده کم است، یک مدل ساده و کمی غلط می‌تواند بهتر از مدل
درست ولی پیچیده پیش‌بینی کند.

\subsubsection{3. تعمیم به‌عنوان معیار}

مدلی را ترجیح می‌دهیم که \textbf{داده‌ی جدید از همان فرایند} را بهتر
پیش‌بینی کند. سه نوع «جدید» با سه ادعای علمی متناظرند:

\begin{longtable}[]{@{}
  >{\raggedright\arraybackslash}p{(\linewidth - 4\tabcolsep) * \real{0.3333}}
  >{\raggedright\arraybackslash}p{(\linewidth - 4\tabcolsep) * \real{0.3333}}
  >{\raggedright\arraybackslash}p{(\linewidth - 4\tabcolsep) * \real{0.3333}}@{}}
\caption{\centering 3. تعمیم به‌عنوان معیار}\tabularnewline
\toprule\noalign{}
\begin{minipage}[b]{\linewidth}\raggedright
داده‌ی جدید یعنی\ensuremath{\ldots}
\end{minipage} & \begin{minipage}[b]{\linewidth}\raggedright
ادعایی که آزمون می‌شود
\end{minipage} & \begin{minipage}[b]{\linewidth}\raggedright
کنار بگذارید\ensuremath{\ldots}
\end{minipage} \\
\midrule\noalign{}
\endfirsthead
\toprule\noalign{}
\begin{minipage}[b]{\linewidth}\raggedright
داده‌ی جدید یعنی\ensuremath{\ldots}
\end{minipage} & \begin{minipage}[b]{\linewidth}\raggedright
ادعایی که آزمون می‌شود
\end{minipage} & \begin{minipage}[b]{\linewidth}\raggedright
کنار بگذارید\ensuremath{\ldots}
\end{minipage} \\
\midrule\noalign{}
\endhead
\bottomrule\noalign{}
\endlastfoot
\lr{trial}های جدید از همان شرکت‌کننده در همان شرایط & مدل فرایند این فرد
را گرفته است & \lr{trial}ها یا آیتم‌های تصادفی \\
\textbf{شرایط} جدید (مثلاً فاصله‌های طولانی‌تر، محرک‌های جدید) & مدل سازوکار
را گرفته، پس می‌تواند برون‌یابی کند & کل شرایط (رویکرد \emph{معیار
تعمیم}) \\
\textbf{شرکت‌کننده‌های} جدید & مدل آنچه را بین آدم‌ها مشترک است گرفته است &
کل شرکت‌کننده‌ها \\
\end{longtable}

\subsubsection{\texorpdfstring{4. \lr{cross-validation
(CV)}}{4. cross-validation (CV)}}

\textbf{اعتبارسنجی \lr{hold-out.}} داده را به یک مجموعه‌ی آموزش (برای
برازش) و یک مجموعه‌ی آزمون (برای ارزیابی) تقسیم کنید. ساده ولی پرهدر و
پرنویز است: نتیجه به یک تقسیم خاص بستگی دارد.

\textbf{اعتبارسنجی \(K\)\lr{-fold.}} داده را به \(K\) بخش \lr{(fold)}
تقریباً برابر تقسیم کنید. برای هر بخش \(k = 1, \dots, K\): 1. مدل را روی
همه‌ی داده‌ها \textbf{جز} بخش \(k\) برازش کنید و
\(\hat{\boldsymbol\theta}_{(-k)}\) را بگیرید. 2. پیش‌بینی‌هایش برای بخش
\(k\) را امتیاز دهید.

امتیازها را روی بخش‌ها جمع کنید. هر مشاهده یک بار برای آزمون استفاده
می‌شود و مدل هرگز روی داده‌ای که با آن آموزش دیده آزمون نمی‌شود. مقادیر
معمول \(K = 5\) یا \(10\) است. \textbf{\lr{Leave-one-out (LOO)}} حالت
خاصی است که هر مشاهده خودش یک بخش است.

\textbf{امتیاز.} برای مدل‌های مبتنی بر درست‌نمایی از \textbf{لگاریتم
درست‌نمایی خارج از نمونه} (امتیاز پیش‌بینی لگاریتمی) استفاده کنید:

\[
\mathrm{CV\text{-}NLL} = -\sum_{k=1}^{K}\ \sum_{i \in \text{\lr{fold }} k} \ln p\big(y_i \mid \hat{\boldsymbol\theta}_{(-k)}, \mathcal{M}\big).
\]

کمتر بهتر است. برای مدل‌های گاوسی، مربع خطای خارج از نمونه گزینه‌ی معادلی
است.

\textbf{قواعد \lr{CV} صادقانه:} 1. \textbf{هر چیزی} که از داده یاد گرفته
می‌شود، از جمله برآورد پارامترها، انتخاب مقدار شروع از شبکه‌ها و تبدیل‌های
داده‌محور، باید \emph{داخل} هر بخش آموزش انجام شود. وگرنه اطلاعات از
مجموعه‌ی آزمون «نشت» می‌کند \lr{(data leakage).} 2. در امتداد واحد درست
تقسیم کنید (بخش 3). برای \lr{CV} در سطح \lr{trial} با شمارش‌های دوجمله‌ای،
\textbf{آیتم‌ها} را داخل هر شرط تقسیم کنید تا هر بخش آموزش هنوز همه‌ی
شرایط را داشته باشد. 3. برآوردهای \lr{CV} پرنویزند. با چند تقسیم تصادفی
تکرار کنید و میانگین و پراکندگی، یا \lr{SE} \emph{تفاوت} مدل‌ها بین
بخش‌ها، را گزارش کنید. 4. با تقسیم در سطح آیتم، \lr{CV} مدل‌هایی را مقایسه
می‌کند که روی \((K-1)/K\) داده برازش شده‌اند، پس مدل‌هایی را که داده‌ی زیادی
لازم دارند کمی جریمه می‌کند.

\begin{quote}
\textbf{پیشرفته: \lr{LOO}، \lr{AIC} و بیز.} برای مدل‌های منظم و \(N\)
بزرگ، \lr{LOO} با امتیاز لگاریتمی و \lr{AIC} (هفته‌ی 7) به‌طور مجانبی یک
مدل را انتخاب می‌کنند. در مدل‌سازی بیزی، تقریبی کارا از \lr{LOO
(PSIS-LOO)} از نمونه‌های پسین محاسبه می‌شود.
\end{quote}

\paragraph{خواندن ریاضی}

\(\mathrm{CV\text{-}NLL} = -\sum_{k}\sum_{i \in \text{\lr{fold }} k} \ln p(y_i \mid \hat{\boldsymbol\theta}_{(-k)}, \mathcal{M})\):

\begin{itemize}
\tightlist
\item
  \(\hat{\boldsymbol\theta}_{(-k)}\): پارامترهایی که \textbf{بدون} بخش
  \(k\) برآورد شده‌اند (علامت منفی در زیرنویس یعنی «کنار گذاشتن»).
\item
  \(\ln p(y_i \mid \hat{\boldsymbol\theta}_{(-k)})\): مدلی که بدون این
  مشاهده‌ها آموزش دیده، آن‌ها را چقدر محتمل می‌داند. مدلی که آن‌ها را خوب
  پیش‌بینی کرده، احتمال بالایی به آن‌ها می‌دهد.
\item
  جمع دوگانه این مقدار را روی همه‌ی بخش‌ها و همه‌ی نقطه‌های کنارگذاشته جمع
  می‌کند؛ علامت منفی آن را به کمیتی برای \emph{کمینه کردن} تبدیل می‌کند،
  مثل \lr{NLL} هفته‌ی 4. همان \lr{NLL} است، ولی هر مشاهده با مدلی سنجیده
  می‌شود که هرگز آن را ندیده.
\end{itemize}

\subsubsection{5. تفسیر یک مقایسه‌ی مدل}

\begin{itemize}
\tightlist
\item
  تفاوت‌های \lr{CV-NLL} بر حسب \textbf{\lr{nat}} هستند. به‌عنوان راهنمای
  تقریبی، تفاوت‌های کمتر از حدود 1 تا 2 \lr{nat}، یا کمتر از دو برابر
  \lr{SE} آن‌ها بین بخش‌ها یا تکرارها، نباید بیش از حد تفسیر شوند.
\item
  مدلی که در \lr{CV} برنده می‌شود بهترین \textbf{پیش‌بینی‌کننده} در میان
  مدل‌های مقایسه‌شده، برای این نوع تعمیم و این مقدار داده، است. «مدل
  واقعی» نیست. با داده‌ی بیشتر، مدل‌های پیچیده‌تر اغلب شروع به بردن می‌کنند،
  چون واریانسشان کم می‌شود.
\item
  \lr{CV} را با بررسی‌های کیفی‌ای که می‌شناسید ترکیب کنید: الگوی
  باقی‌مانده‌ها، معقول بودن پارامترها، و قابل تفسیر بودن مدل.
\end{itemize}

\paragraph{برای یک مهندس}

\begin{itemize}
\tightlist
\item
  \textbf{\lr{overfitting}} همان مشکل آشنای برازش یک چندجمله‌ای با درجه‌ی
  بالا یا یک فیلتر با تعداد \lr{tap} زیاد روی یک سیگنال کوتاه و پرنویز
  است: خطای آموزش به صفر می‌رسد و خطای پیش‌بینی منفجر می‌شود.
\item
  \textbf{\lr{trade-off} بین \lr{bias} و \lr{variance}} همان چیزی است که
  هنگام انتخاب مرتبه‌ی مدل (مرتبه‌ی \lr{AR}، تعداد \lr{tap}های فیلتر) یا
  پنجره‌ی هموارسازی می‌بینید: درجه‌ی آزادی کم \lr{bias} می‌دهد و درجه‌ی آزادی
  زیاد \lr{variance.}
\item
  \textbf{\lr{cross-validation}} نسخه‌ی داده‌محور انتخاب مرتبه‌ی مدل است.
  معیارهایی از نوع \lr{AIC/MDL} (هفته‌ی 7) میان‌برهای تحلیلی‌اند که شاید از
  انتخاب مرتبه‌ی مدل \lr{AR} بشناسید.
\item
  \textbf{آموزش روی داده‌ی آزمون} مثل این است که یک گیرنده را روی همان
  تحقق نویزی تنظیم کنید که بعداً \lr{BER} آن را رویش اندازه می‌گیرید:
  کارایی گزارش‌شده خوش‌بینانه است.
\item
  \textbf{\lr{nat}} نسخه‌ی لگاریتم طبیعی \lr{bit} است: تفاوت 1 \lr{nat}
  یعنی نسبت درست‌نمایی \(e \approx 2.7\).
\end{itemize}

\paragraph{برای مدرس}

\begin{itemize}
\tightlist
\item
  \textbf{با مثال حل‌شده‌ی چندجمله‌ای روی تخته شروع کنید.} چهار نقطه، چهار
  مدل. بگذارید دانشجوها بهترین مدل را با \lr{SSE} آموزش انتخاب کنند
  (مکعبی را انتخاب می‌کنند)، بعد بپرسید پیش‌بینی آن در 16 روز چیست
  (\(-8.5\)). خنده‌ی کلاس همان درس است.
\item
  \textbf{\lr{leave-one-out} را برای یک نقطه با دست انجام دهید} (نقطه‌ی 8
  روز را با دست بپوشانید، خط را با چشم دوباره برازش کنید و پیش‌بینی
  کنید).
\item
  \textbf{روی «چه چیزی کنار گذاشته می‌شود» تأکید کنید.} سه نوع داده‌ی
  «جدید» (بخش 3) را روی تخته بنویسید؛ بیشتر اشتباه‌های پروژه‌های دانشجویی
  از کنار گذاشتن چیز اشتباه می‌آید.
\item
  \textbf{سؤال‌های رایج.}

  \begin{itemize}
  \tightlist
  \item
    \emph{«اگر مدل درست کمتر \lr{overfit} می‌کند، چرا یک مدل غلط می‌تواند
    در \lr{CV} برنده شود؟»} با داده‌ی کم، یک مدل ساده‌ی غلط ممکن است
    \lr{variance} کمتری داشته باشد. \lr{CV} بهترین \emph{پیش‌بینی‌کننده}
    را انتخاب می‌کند، نه مدل درست را.
  \item
    \emph{«چند بخش \lr{(fold)}؟»} 5 یا 10 استاندارد است؛ برای داده‌های
    خیلی کوچک \lr{LOO.}
  \item
    \emph{«آیا \(R^2 = 0.98\) خوب است؟»} برای یک منحنی نزولی، تقریباً هر
    مدل یکنوایی \(R^2\) بالا می‌گیرد؛ چیز زیادی نمی‌گوید.
  \end{itemize}
\end{itemize}

\begin{center}\rule{0.5\linewidth}{0.5pt}\end{center}

\subsection{مثال حل‌شده}

\textbf{داده.} چهار نسبت یادآوری \(y = (0.62, 0.50, 0.47, 0.27)\) در
\(t = (1, 2, 4, 8)\) روز. چندجمله‌ای‌هایی با درجه‌ی رو به افزایش را به‌عنوان
«مدل» مقایسه می‌کنیم.

\subsubsection{قدم 1: برازش روی همه‌ی داده‌ها}

\begin{longtable}[]{@{}llll@{}}
\caption{\centering قدم 1: برازش روی همه‌ی داده‌ها}\tabularnewline
\toprule\noalign{}
مدل & \(k\) & معادله‌ی برازش‌شده & \lr{SSE} آموزش \\
\midrule\noalign{}
\endfirsthead
\toprule\noalign{}
مدل & \(k\) & معادله‌ی برازش‌شده & \lr{SSE} آموزش \\
\midrule\noalign{}
\endhead
\bottomrule\noalign{}
\endlastfoot
ثابت & 1 & \(\hat y = 0.465\) & 0.0633 \\
خطی & 2 & \(\hat y = 0.637 - 0.046\,t\) & 0.0032 \\
درجه‌ی دو & 3 & \(\hat y = 0.647 - 0.053\,t + 0.0008\,t^2\) & 0.0031 \\
درجه‌ی سه & 4 & دقیقاً از هر 4 نقطه می‌گذرد & \textbf{0.0000} \\
\end{longtable}

بر اساس \lr{SSE} آموزش، چندجمله‌ای درجه‌ی سه «کامل» است. ولی از آن بخواهید
یادآوری را در \(t = 16\) روز پیش‌بینی کند: \(-8.5\) پیش‌بینی می‌کند، یعنی
\emph{منفی 850\% یادآوری}. مدل خطی \(-0.10\) پیش‌بینی می‌کند (آن هم ناممکن
ولی کمتر پوچ) و مدل ثابت 0.47.

\subsubsection{\texorpdfstring{قدم 2: اعتبارسنجی
\lr{leave-one-out}}{قدم 2: اعتبارسنجی leave-one-out}}

هر بار یک نقطه را کنار بگذارید، روی سه نقطه‌ی باقی‌مانده دوباره برازش کنید
و نقطه‌ی کنارگذاشته را پیش‌بینی کنید.

\begin{longtable}[]{@{}
  >{\raggedright\arraybackslash}p{(\linewidth - 8\tabcolsep) * \real{0.2000}}
  >{\raggedright\arraybackslash}p{(\linewidth - 8\tabcolsep) * \real{0.2000}}
  >{\raggedright\arraybackslash}p{(\linewidth - 8\tabcolsep) * \real{0.2000}}
  >{\raggedright\arraybackslash}p{(\linewidth - 8\tabcolsep) * \real{0.2000}}
  >{\raggedright\arraybackslash}p{(\linewidth - 8\tabcolsep) * \real{0.2000}}@{}}
\caption{\centering قدم 2: اعتبارسنجی \lr{leave-one-out}}\tabularnewline
\toprule\noalign{}
\begin{minipage}[b]{\linewidth}\raggedright
\(t\) کنارگذاشته
\end{minipage} & \begin{minipage}[b]{\linewidth}\raggedright
مشاهده‌شده
\end{minipage} & \begin{minipage}[b]{\linewidth}\raggedright
پیش‌بینی ثابت
\end{minipage} & \begin{minipage}[b]{\linewidth}\raggedright
پیش‌بینی خطی
\end{minipage} & \begin{minipage}[b]{\linewidth}\raggedright
پیش‌بینی درجه‌ی دو
\end{minipage} \\
\midrule\noalign{}
\endfirsthead
\toprule\noalign{}
\begin{minipage}[b]{\linewidth}\raggedright
\(t\) کنارگذاشته
\end{minipage} & \begin{minipage}[b]{\linewidth}\raggedright
مشاهده‌شده
\end{minipage} & \begin{minipage}[b]{\linewidth}\raggedright
پیش‌بینی ثابت
\end{minipage} & \begin{minipage}[b]{\linewidth}\raggedright
پیش‌بینی خطی
\end{minipage} & \begin{minipage}[b]{\linewidth}\raggedright
پیش‌بینی درجه‌ی دو
\end{minipage} \\
\midrule\noalign{}
\endhead
\bottomrule\noalign{}
\endlastfoot
1 & 0.62 & 0.413 & 0.560 & 0.498 \\
2 & 0.50 & 0.453 & 0.570 & 0.570 \\
4 & 0.47 & 0.463 & 0.448 & 0.330 \\
8 & 0.27 & 0.530 & 0.275 & \textbf{1.250} \\
\textbf{\lr{SSE} در \lr{LOO}} & & \textbf{0.1125} & \textbf{0.0090} &
\textbf{1.0000} \\
\end{longtable}

(چندجمله‌ای درجه‌ی سه اینجا قابل اعتبارسنجی نیست: سه نقطه چهار پارامتر را
تعیین نمی‌کنند.)

\textbf{تفسیر.} \lr{SSE} آموزش مدل درجه‌ی دو عملاً با مدل خطی برابر است،
ولی خطای \lr{LOO} آن بیش از 100 برابر بزرگ‌تر است. وقتی نقطه‌ی 8 روزه را
نمی‌بیند، انحنایش را طوری خم می‌کند که سه فاصله‌ی کوتاه را برازش کند و در 8
روز یادآوری \emph{بیشتر از 1} پیش‌بینی می‌کند. پارامتر اضافه نویز را برازش
کرد، نه فراموشی را. مدل خطی از میان این سه بهترین تعمیم را دارد. دقیقاً
به همین دلیل به‌جایشان از مدل‌های نمایی و توانی استفاده می‌کنیم.)

\begin{center}\rule{0.5\linewidth}{0.5pt}\end{center}

\Needspace{11\baselineskip}

\subsection{\texorpdfstring{\lr{lab} (2 ساعت)}{lab (2 ساعت)}}

\subsubsection{اهداف}

\begin{enumerate}
\def\labelenumi{\arabic{enumi}.}
\tightlist
\item
  دیدن \lr{overfitting} در عمل: خطای آموزش در برابر خطای آزمون با افزایش
  درجه‌ی چندجمله‌ای.
\item
  صورت‌بندی \lr{CV} \(K\)\lr{-fold} در سطح آیتم با لگاریتم درست‌نمایی
  دوجمله‌ای خارج از نمونه، و مقایسه‌ی چهار مدل فراموشی.
\item
  اجرای یک آزمون \textbf{معیار تعمیم}: آموزش روی فاصله‌های کوتاه، پیش‌بینی
  فاصله‌های طولانی.
\end{enumerate}

\textbf{زمان‌بندی:} بخش \lr{A} 25 دقیقه · بخش \lr{B} 45 دقیقه · بخش
\lr{C} 25 دقیقه · بخش \lr{D} 25 دقیقه.

\subsubsection{آماده‌سازی: داده و چهار مدل}

بر اساس \lr{NLL} آموزش، مدل‌های پیچیده‌تر برنده‌اند. دقت کنید که بهترین
برازش مدل دوردی یک جزء کند با نرخ 0 دارد: مثل هفته‌ی 3 به مدل نمایی با
مجانب تبدیل شده است (با \lr{NLL} یکسان). حالا ببینیم آیا \emph{پیش‌بینی}
بهتری هم دارند.

\subsubsection{\texorpdfstring{بخش \lr{A: overfitting} با چندجمله‌ای‌ها
(25 دقیقه)}{بخش A: overfitting با چندجمله‌ای‌ها (25 دقیقه)}}

چندجمله‌ای‌هایی بر حسب \(x = \ln(1+t)\) را به‌عنوان مدل‌های منعطفِ عمداً
«بی‌فکر» به کار می‌بریم. داده‌های \textbf{آزمون} 500 مجموعه‌داده‌ی تازه‌اند که
با همان طرح از مدل توانی واقعی شبیه‌سازی شده‌اند.

\textbf{به این دقت کنید:} خطای آموزش با هر پارامتر اضافه کم می‌شود و در
\(k = 8\) به 0 می‌رسد (درون‌یابی). خطای آزمون در یک پیچیدگی میانی کمترین
است و بعد بالا می‌رود: پارامترهای اضافه دنبال نویز می‌دوند. حتی بهترین
چندجمله‌ای هم بدتر از مدل توانی \emph{واقعی} پیش‌بینی می‌کند؛ \lr{MSE}
آزمون آن را می‌توانید به‌عنوان تمرین حساب کنید.

\subsubsection{\texorpdfstring{بخش \lr{B: cross-validation}
\(K\)\lr{-fold} روی آیتم‌ها (45
دقیقه)}{بخش B: cross-validation K-fold روی آیتم‌ها (45 دقیقه)}}

در هر فاصله 50 آیتم آزمون شده است. هر آیتم را \emph{داخل} هر فاصله به
یکی از \(K = 5\) بخش اختصاص می‌دهیم تا هر مجموعه‌ی آموزش همه‌ی 8 فاصله را
داشته باشد.

دقت کنید که \lr{NLL} آموزش (دوجمله‌ای، با ضرایب دوجمله‌ای) و \lr{CV-NLL}
(برنولی، در سطح آیتم) مقیاس‌های متفاوتی دارند. مدل‌ها را \textbf{داخل} هر
ستون مقایسه کنید، نه بین ستون‌ها.

\textbf{تفسیر:} 1. آیا برتری‌اش بر نفر دوم از تغییرپذیری بین تکرارهای
\lr{CV} بیشتر است؟ 2. مدل دوردی بهترین برازش آموزش را دارد. در \lr{CV}
چه بر سرش می‌آید و چرا؟ 3. آیا \lr{CV} با این داده‌ها می‌تواند مدل توانی و
مدل نمایی با مجانب را از هم جدا کند؟ چه چیزی کمک می‌کرد؟

\subsubsection{\texorpdfstring{بخش \lr{C:} معیار تعمیم: پیش‌بینی فاصله‌های
طولانی‌تر (25
دقیقه)}{بخش C: معیار تعمیم: پیش‌بینی فاصله‌های طولانی‌تر (25 دقیقه)}}

مدلی که فهمیده فراموشی \emph{چگونه} کار می‌کند باید بتواند برون‌یابی کند.
روی پنج فاصله‌ی کوتاه‌تر (\(t \le 7\) روز) آموزش دهید و سه فاصله‌ی طولانی‌تر
(14، 21، 30 روز) را که مدل‌ها هرگز ندیده‌اند پیش‌بینی کنید.

\textbf{به این دقت کنید:} اینجا مدل \textbf{نمایی با مجانب} بهترین
برون‌یابی را دارد، \emph{نه} مدل توانی واقعی. سه نقطه‌ی کنارگذاشته اتفاقاً
صاف‌اند (0.26، 0.26، 0.28) و منحنی‌ای که ثابت می‌شود آن‌ها را خوب پیش‌بینی
می‌کند. مدل توانی، مثل فرایند واقعی، به افت ادامه می‌دهد، ولی در این نمونه
نویز فاصله‌های طولانی یادآوری را صاف نشان داده است. مدل نمایی خالص و مدل
دوردی (که برازش فاصله‌های کوتاهش هیچ قیدی روی جزء کند نمی‌گذارد) بد
برون‌یابی می‌کنند. درس‌ها: (1) معیارهای تعمیم مختلف می‌توانند مدل‌ها را
متفاوت رتبه‌بندی کنند؛ (2) یک آزمون روی سه نقطه‌ی داده خودش پرنویز است، پس
برد یک مدل در یک مجموعه‌ی آزمون اثبات نیست؛ (3) توافق چند معیار (بخش
\lr{B}، بخش \lr{C}، و \lr{AIC/BIC} هفته‌ی 7) قانع‌کننده‌تر از هر یک از
آن‌هاست.

\subsubsection{\texorpdfstring{بخش \lr{D:} کنار هم گذاشتن (25
دقیقه)}{بخش D: کنار هم گذاشتن (25 دقیقه)}}

\begin{landscape}\small
\begin{longtable}[]{@{}
  >{\raggedright\arraybackslash}p{(\linewidth - 12\tabcolsep) * \real{0.1429}}
  >{\raggedright\arraybackslash}p{(\linewidth - 12\tabcolsep) * \real{0.1429}}
  >{\raggedright\arraybackslash}p{(\linewidth - 12\tabcolsep) * \real{0.1429}}
  >{\raggedright\arraybackslash}p{(\linewidth - 12\tabcolsep) * \real{0.1429}}
  >{\raggedright\arraybackslash}p{(\linewidth - 12\tabcolsep) * \real{0.1429}}
  >{\raggedright\arraybackslash}p{(\linewidth - 12\tabcolsep) * \real{0.1429}}
  >{\raggedright\arraybackslash}p{(\linewidth - 12\tabcolsep) * \real{0.1429}}@{}}
\caption{\centering بخش \lr{D: }کنار هم گذاشتن (25 دقیقه)}\tabularnewline
\toprule\noalign{}
\begin{minipage}[b]{\linewidth}\raggedright
مدل
\end{minipage} & \begin{minipage}[b]{\linewidth}\raggedright
\(k\)
\end{minipage} & \begin{minipage}[b]{\linewidth}\raggedright
\lr{NLL} آموزش
\end{minipage} & \begin{minipage}[b]{\linewidth}\raggedright
\lr{CV-NLL} پنج‌بخشی (میانگین \ensuremath{\pm} \lr{SD)}
\end{minipage} & \begin{minipage}[b]{\linewidth}\raggedright
\lr{NLL} آزمون فاصله‌های طولانی
\end{minipage} & \begin{minipage}[b]{\linewidth}\raggedright
الگوی باقی‌مانده‌ها (هفته‌ی 4)
\end{minipage} & \begin{minipage}[b]{\linewidth}\raggedright
پارامترها معقول‌اند؟
\end{minipage} \\
\midrule\noalign{}
\endfirsthead
\toprule\noalign{}
\begin{minipage}[b]{\linewidth}\raggedright
مدل
\end{minipage} & \begin{minipage}[b]{\linewidth}\raggedright
\(k\)
\end{minipage} & \begin{minipage}[b]{\linewidth}\raggedright
\lr{NLL} آموزش
\end{minipage} & \begin{minipage}[b]{\linewidth}\raggedright
\lr{CV-NLL} پنج‌بخشی (میانگین \ensuremath{\pm} \lr{SD)}
\end{minipage} & \begin{minipage}[b]{\linewidth}\raggedright
\lr{NLL} آزمون فاصله‌های طولانی
\end{minipage} & \begin{minipage}[b]{\linewidth}\raggedright
الگوی باقی‌مانده‌ها (هفته‌ی 4)
\end{minipage} & \begin{minipage}[b]{\linewidth}\raggedright
پارامترها معقول‌اند؟
\end{minipage} \\
\midrule\noalign{}
\endhead
\bottomrule\noalign{}
\endlastfoot
\lr{exp} & & & & & & \\
\lr{pow} & & & & & & \\
\lr{expasym} & & & & & & \\
\lr{two} & & & & & & \\
\end{longtable}
\end{landscape}

\subsubsection{تمرین‌ها (به ترتیب سختی)}

\textbf{تمرین 1 \ensuremath{\star}.} در بخش \lr{A}، مدل توانی واقعی با پارامترهای واقعی
را به‌عنوان یک «مدل» اضافه کنید (هیچ پارامتر برازش‌شده‌ای ندارد). \lr{MSE}
آزمون آن را حساب کنید و با بهترین چندجمله‌ای مقایسه کنید. خطای کاهش‌ناپذیر
چقدر است؟

\textbf{تمرین 2 \ensuremath{\star}.} بخش \lr{B} را با \(K = 2\)، \(K = 10\) و
\lr{leave-one-item-out} (\(K = 50\)) دوباره اجرا کنید. آیا رتبه‌بندی
تغییر می‌کند؟

\textbf{تمرین 3 \ensuremath{\star}\ensuremath{\star} (نشت داده).} یک \lr{CV} «غلط» بنویسید که در آن مقادیر
شروع هر بخش از برازش روی \emph{کل} مجموعه‌داده گرفته شود، و یکی دیگر که
در آن آیتم‌های آزمون هر بخش اشتباهاً در آموزش هم باشند. نتایج را با
\lr{CV} درست مقایسه کنید.

\textbf{تمرین 4 \ensuremath{\star}\ensuremath{\star} (داده‌ی بیشتر).} یک مجموعه‌داده‌ی جدید از مدل توانی
واقعی با \(n = 500\) آیتم در هر فاصله شبیه‌سازی کنید و بخش‌های \lr{B} و
\lr{C} را دوباره اجرا کنید. آیا رتبه‌بندی روشن‌تر می‌شود؟ آیا مدل دوردی
هیچ‌وقت برنده می‌شود؟

\textbf{تمرین 5 \ensuremath{\star}\ensuremath{\star}\ensuremath{\star} (پیشرفته: \lr{leave-one-participant-out).}} از مولد
داده‌ی تمرین بخش \lr{D} هفته‌ی 5 با 20 شرکت‌کننده استفاده کنید. \lr{CV} ای
صورت‌بندی کنید که هر بخش آن \emph{یک شرکت‌کننده} را کنار بگذارد: یک مدل
تجمیعی (میانگین گروه) را روی 19 شرکت‌کننده برازش کنید و پیش‌بینی منحنی
شرکت‌کننده‌ی کنارگذاشته را با لگاریتم درست‌نمایی گاوسی امتیاز دهید. مدل‌های
نمایی و توانی را مقایسه کنید. بعد همین کار را با برازش‌های سطح فرد انجام
دهید. توضیح دهید چرا این دو طرح \lr{CV} به دو سؤال علمی متفاوت جواب
می‌دهند.

\subsubsection{\texorpdfstring{\lr{lab check-off}}{lab check-off}}

نشان دهید: (1) نمودار آموزش در برابر آزمون از بخش \lr{A}، (2) جدول
خلاصه‌ی \lr{CV} از بخش \lr{B}، و (3) پاراگراف بخش \lr{D.}

\begin{center}\rule{0.5\linewidth}{0.5pt}\end{center}

\subsection{\texorpdfstring{ارتباط با هوش مصنوعی و
\lr{LLM}ها}{ارتباط با هوش مصنوعی و LLMها}}

\textbf{آلودگی \lr{benchmark (contamination)} همان نشت داده است.}
\lr{LLM}ها روی حجم عظیمی از متن اینترنت آموزش می‌بینند. اگر سؤال‌های یک
\lr{benchmark} (یا جواب‌هایشان) آنلاین منتشر شده باشند، ممکن است در داده‌ی
آموزشی باشند، پس مدل روی آیتم‌هایی آزموده می‌شود که قبلاً دیده. این دقیقاً
مسئله‌ی نشت در \lr{cross-validation} است: مجموعه‌ی آزمون دیگر از مجموعه‌ی
آموزش مستقل نیست، پس نمره علاوه بر تعمیم، حفظ کردن را هم می‌سنجد. این یکی
از جدی‌ترین مشکلات ارزیابی سیستم‌های هوش مصنوعی است، از جمله ادعاهایی که
می‌گویند \lr{LLM}ها روی تکلیف‌های کلاسیک روان‌شناسی، که مواد آن‌ها همه‌جا
منتشر شده، توانایی‌های شناختی انسان‌گونه نشان می‌دهند.

نمره‌ی گزارش‌شده (0.73) توانایی حل مسئله‌های \emph{تازه} (0.63) را بیش از
واقع نشان می‌دهد. راه‌حل همان است که در \lr{lab} این هفته دیدید: ارزیابی
روی آیتم‌هایی که مطمئنیم در برازش استفاده نشده‌اند، مثلاً آیتم‌های
تازه‌نوشته، یا گونه‌های نظام‌مند تکلیف‌های کلاسیک.

آیا فقط از روی دقت گزارش‌شده می‌توانید بگویید چقدر نشت کرده؟ فرض کنید برای
هر آیتم یک گونه‌ی «تازه» (همان منطق، عبارت‌بندی تازه) داریم. گونه‌های تازه
را با \(P(\text{\lr{correct}}) = 0.60\) شبیه‌سازی کنید و اثر نشت را به‌صورت
تفاوت دقت آیتم‌های اصلی و تازه، با یک بازه‌ی \lr{bootstrap} جفتی (هفته‌ی 5)
برآورد کنید. - \ensuremath{\star}\ensuremath{\star} بحث: یک \lr{LLM} معمای استدلالی معروفی از یک مقاله‌ی
روان‌شناسی دهه‌ی 1970 را حل می‌کند. چرا این شاهد ضعیفی است بر اینکه «مثل
انسان استدلال می‌کند»؟ چه آزمونی قوی‌تر است؟

\begin{center}\rule{0.5\linewidth}{0.5pt}\end{center}

\subsection{اشتباهات رایج و برداشت‌های نادرست}

\begin{enumerate}
\def\labelenumi{\arabic{enumi}.}
\tightlist
\item
  \textbf{«مدلی که بهترین برازش را دارد بهترین مدل است.»} برازش انعطاف
  را پاداش می‌دهد. همیشه بپرسید مدل داده‌ی جدید را چقدر خوب پیش‌بینی می‌کند.
\item
  \textbf{«انعطاف = تعداد پارامترها.»} شکل تابعی هم مهم است. مدل‌هایی با
  \(k\) برابر می‌توانند در آنچه می‌توانند برازش کنند خیلی فرق داشته باشند.
\item
  \textbf{نشت اطلاعات آزمون به آموزش،} مثلاً انتخاب مقدار شروع، کران یا
  تبدیل از روی کل داده، یا نرمال‌سازی با آماره‌های کل داده.
\item
  \textbf{تقسیم در امتداد واحد غلط.} اگر ادعا درباره‌ی شرکت‌کننده‌های جدید
  است، کنار گذاشتن \lr{trial}های تصادفی تعمیم را بیش از حد برآورد می‌کند.
\item
  \textbf{بیش‌تفسیر تفاوت‌های کوچک \lr{CV.}} برآوردهای \lr{CV} پرنویزند.
  تقسیم‌ها را تکرار کنید و به تغییرپذیری تفاوت‌ها نگاه کنید.
\item
  \textbf{\lr{«CV} مدل واقعی را انتخاب کرد.»} \lr{CV} بهترین
  \emph{پیش‌بینی‌کننده} را برای این مقدار داده انتخاب می‌کند. با داده‌ی کم،
  یک مدل ساده‌تر (حتی غلط) می‌تواند بهتر از مدل واقعی پیش‌بینی کند.
\item
  \textbf{استفاده از \(R^2\) برای مقایسه‌ی مدل‌های شناختی غیرخطی.} بین
  مدل‌ها قابل مقایسه نیست و انعطاف را پاداش می‌دهد.
\end{enumerate}

\begin{center}\rule{0.5\linewidth}{0.5pt}\end{center}

\subsection{خودآزمایی}

\textbf{1.} چرا وقتی به یک مدل \lr{nested} پارامتر اضافه می‌کنیم، خطای
آموزش همیشه کم می‌شود (یا ثابت می‌ماند)؟ \textbf{پاسخ} مدل بزرگ‌تر مدل
کوچک‌تر را به‌عنوان حالت خاص در بر دارد (مثلاً با صفر گذاشتن پارامتر
اضافه)، پس بهترین برازشش هرگز بدتر نیست. بهینه‌ساز همیشه می‌تواند برازش
مدل کوچک‌تر را بازتولید کند و معمولاً با برازش نویز آن را بهتر می‌کند.

\textbf{2.} \lr{overfitting} را بر حسب علامت و نویز تعریف کنید.
\textbf{پاسخ} داده = علامت (فرایند نظام‌مند) + نویز (تغییرات تصادفی).
مدلی که \lr{overfitting} می‌کند خودش را با نویز نمونه‌ی خاص تطبیق می‌دهد.
این برازشش روی این داده‌ها را بهتر و پیش‌بینی داده‌ی جدید را، که نویزش
متفاوت است، بدتر می‌کند.

\textbf{3.} \lr{cross-validation} \(K\)\lr{-fold} را در چهار قدم توضیح
دهید. \textbf{پاسخ} (1) داده را به \(K\) بخش تقسیم کنید. (2) برای هر
بخش، مدل را روی \(K-1\) بخش دیگر برازش کنید. (3) پیش‌بینی‌های مدل برازش‌شده
را برای بخش کنارگذاشته امتیاز دهید (مثلاً با لگاریتم درست‌نمایی). (4)
امتیازها را روی بخش‌ها جمع یا میانگین کنید؛ در صورت امکان با تقسیم‌های
تصادفی جدید تکرار کنید.

\textbf{4.} در آزمایش فراموشی، چرا آیتم‌ها را \emph{داخل} هر فاصله‌ی
نگهداری به بخش‌ها اختصاص می‌دهیم به‌جای اینکه کل فاصله‌ها را کنار بگذاریم؟
\textbf{پاسخ} تقسیم داخل فاصله پیش‌بینی \lr{trial}های جدید در همان شرایط
را آزمون می‌کند، و هر بخش آموزش هنوز همه‌ی شرایط را دارد. کنار گذاشتن کل
فاصله‌ها ادعای دیگری را آزمون می‌کند: برون‌یابی به شرایط جدید (معیار تعمیم،
بخش \lr{C).}

\textbf{5.} \lr{CV-NLL} مدل \lr{A} به‌طور میانگین 2.1 \lr{nat} کمتر از
مدل \lr{B} است، ولی این تفاوت بین تکرارهای \lr{CV} با \lr{SD} برابر 3.0
تغییر می‌کند. چه نتیجه‌ای می‌گیرید؟ \textbf{پاسخ} شاهد ضعیف است. به‌طور
دقیق، \lr{SD} بین تکرارها تغییرپذیری بین تقسیم‌ها را نشان می‌دهد نه عدم
قطعیت نمونه‌گیری داده را، ولی در هر حال یک تفاوت 2 \lr{nat} با این همه
ناپایداری نباید بیش‌تفسیر شود. داده‌ی بیشتری جمع کنید یا مدل‌ها را تقریباً
هم‌ارز در پیش‌بینی گزارش کنید.

\textbf{6.} چرا یک مدل ساده‌تر و «غلط» گاهی در \lr{CV} از مدل واقعی و
پیچیده‌تر بهتر عمل می‌کند؟ \textbf{پاسخ} با داده‌ی محدود، پارامترهای مدل
پیچیده با واریانس بالا برآورد می‌شوند و این واریانس بیشتر از سوگیری مدل
ساده‌تر به پیش‌بینی آسیب می‌زند \lr{(trade-off} سوگیری\lr{--}واریانس). با داده‌ی
بیشتر واریانس مدل واقعی کم می‌شود و بالاخره برنده می‌شود.

\textbf{7.} دانشجویی \(R^2 = 0.97\) را برای یک مدل غیرخطی فراموشی گزارش
می‌کند و نتیجه می‌گیرد مدل درست است. دو اعتراض بیاورید. \textbf{پاسخ} (1)
\(R^2\) برای تقریباً هر منحنی نزولی یکنوایی بالاست، چون بیشترِ واریانس
داده‌های فراموشی همان افت کلی است. (2) برازش خوب چیزی درباره‌ی تعمیم یا
انعطاف نمی‌گوید؛ الگوی باقی‌مانده‌ها، \lr{CV} و مقایسه با مدل‌های جایگزین
لازم است.

\begin{center}\rule{0.5\linewidth}{0.5pt}\end{center}

\Needspace{11\baselineskip}

\subsection{تکلیف}

\subsubsection{\texorpdfstring{\lr{A3:} مقایسه و بازیابی مدل (1.0 نمره
از 20؛ انتشار هفته‌ی 6، موعد ابتدای درس هفته‌ی
9)}{A3: مقایسه و بازیابی مدل (1.0 نمره از 20؛ انتشار هفته‌ی 6، موعد ابتدای درس هفته‌ی 9)}}

\lr{A3} این هفته و هفته‌ی بعد را در بر می‌گیرد: کارهای 1 تا 3 از مطالب این
هفته و کارهای 4 تا 6 از مطالب هفته‌ی 7 \lr{(AIC/BIC}، بازیابی پارامتر و
مدل) استفاده می‌کنند.

\textbf{کارها.}

\begin{enumerate}
\def\labelenumi{\arabic{enumi}.}
\tightlist
\item
  \textbf{برازش‌های فردی (1.5 امتیاز).} مدل‌های نمایی، توانی و نمایی با
  مجانب را برای هر شرکت‌کننده با \lr{MLE} چندشروعی برازش کنید. جدولی از
  پارامترها و \lr{NLL}ها گزارش کنید؛ برآوردهای روی کران را مشخص کنید.
\item
  \textbf{برازش مطلق (1 امتیاز).} \(G^2\) را برای هر شرکت‌کننده و مدل
  حساب کنید. هر مدل برای چند شرکت‌کننده برازش ناموفق دارد (\(p < .05\))؟
\item
  \textbf{\lr{cross-validation} (2 امتیاز).} برای هر شرکت‌کننده، \lr{CV}
  پنج‌بخشی در سطح آیتم (5 بار تکرار) را برای سه مدل اجرا کنید. گزارش کنید
  هر مدل برای چند شرکت‌کننده برنده است، و مجموع \lr{CV-NLL} روی
  شرکت‌کننده‌ها را با تغییرپذیری‌اش.
\item
  \textbf{معیارهای اطلاعاتی (1.5 امتیاز، هفته‌ی 7).} \lr{AIC}، \lr{AICc}
  و \lr{BIC} را برای هر شرکت‌کننده و مدل و وزن‌های آکائیکه را حساب کنید.
  آیا معیارها با \lr{CV} توافق دارند؟
\item
  \textbf{بازیابی (2.5 امتیاز، هفته‌ی 7).} (الف) بازیابی پارامتر برای مدل
  توانی در این طرح: 200 شبیه‌سازی با پارامترهایی از دامنه‌های معقول؛
  مقادیر واقعی در برابر بازیابی‌شده را رسم کنید، همبستگی‌ها را گزارش کنید
  و درباره‌ی هر \lr{trade-off}ی توضیح دهید. (ب) بازیابی مدل برای سه مدل
  با \lr{AIC} و \lr{BIC:} یک \lr{confusion matrix} 3 × 3 از 100 شبیه‌سازی
  برای هر مدل تولیدکننده.
\end{enumerate}

\textbf{روبریک (10 امتیاز، معادل 1.0):} امتیاز هر کار در بالا آمده است.
داخل هر کار: درستی 60\%، وضوح جدول‌ها و نمودارها 20\%، تفسیر 20\%. برای
تحویل با تأخیر، روزی 10\% کسر می‌شود (حداکثر 3 روز).

\begin{center}\rule{0.5\linewidth}{0.5pt}\end{center}

\subsection{مطالعه‌ی بیشتر}

\begin{itemize}
\tightlist
\item
  \textbf{\lr{Farrell} \lr{\&} \lr{Lewandowsky} (2018)، فصل 10.} مقایسه‌ی
  مدل‌ها: چرا برازش کافی نیست، انعطاف‌پذیری، و \lr{cross-validation} در
  کنار دیگر روش‌های مقایسه. \emph{(منبع اصلی.)}
\item
  \textbf{عظیمی‌فر و احمدی (1396)}، جلد 1: مطالب مربوط به تعمیم‌پذیری،
  \lr{overfitting} و \lr{trade-off} سوگیری\lr{--}واریانس از دیدگاه شناسایی
  الگو {[}شماره‌ی فصل بررسی شود{]}.
\item
  \emph{کلاسیک‌های اختیاری، خارج از فهرست درس:} \lr{Roberts, S.,} \lr{\&}
  \lr{Pashler, H.} (2000). \lr{How persuasive is a good fit? A comment
  on theory testing.} \emph{\lr{Psychological Review,} 107}, 358\lr{--}367.
  \lr{Busemeyer, J. R.,} \lr{\&} \lr{Wang, Y.-M.} (2000). \lr{Model
  comparisons and model selections based on generalization criterion
  methodology.} \emph{\lr{Journal of Mathematical Psychology,} 44},
  171\lr{--}189.
\end{itemize}

\emph{پیش‌نمایش هفته‌ی 7:} \lr{cross-validation} کلی است ولی از نظر
محاسباتی پرهزینه. معیارهای اطلاعاتی \lr{(AIC}، \lr{BIC)} تعمیم را از یک
برازش واحد تخمین می‌زنند. و پیش از اعتماد به هر مقایسه‌ای، باید بررسی کنیم
آیا روش ما می‌تواند پارامترها و مدل‌هایی را که واقعاً داده را تولید کرده‌اند
بازیابی کند.
